\section{Manhattan para Chebyshev}

\textbf{Problema base.} A distância de Manhattan entre $(x_1, y_1)$ e $(x_2, y_2)$ é $|x_1 - x_2| + |y_1 - y_2|$. A soma
junta as duas coordenadas, então não dá para tratar $x$ e $y$ separadamente (ex.: achar o par mais distante sem testar
todos os pares).

\textbf{Ideia.} Gire o plano em $45^\circ$ com $u = x + y$ e $v = x - y$. Como $|a| + |b| = \max(|a + b|, |a - b|)$,
Manhattan vira Chebyshev, que é a maior das diferenças:
\[
\boxed{|x_1 - x_2| + |y_1 - y_2| = \max(|u_1 - u_2|,\ |v_1 - v_2|)}
\]
Com o $\max$, cada eixo ($u$ ou $v$) pode ser analisado sozinho.

\subsection*{Variações}

\begin{itemize}
  \item \textbf{Volta (Chebyshev para Manhattan):} com o mesmo $u, v$,
        \[ 2 \max(|x_1 - x_2|,\ |y_1 - y_2|) = |u_1 - u_2| + |v_1 - v_2|. \]
  \item \textbf{Paridade:} $u$ e $v$ têm sempre a mesma paridade; ao voltar com $x = \frac{u+v}{2}$, $y = \frac{u-v}{2}$,
        um $(u, v)$ de paridades diferentes não vem de ponto inteiro.
\end{itemize}

\subsection*{Aplicações}

\begin{itemize}
  \item \textbf{Par mais distante} entre $N$ pontos em $O(N)$: $\max(\max u - \min u,\ \max v - \min v)$.
  \item \textbf{Pontos a distância $\le d$ de $(x_0, y_0)$:} o losango vira o quadrado $|u - u_0| \le d$, $|v - v_0| \le d$
        (com $u_0 = x_0 + y_0$, $v_0 = x_0 - y_0$), alinhado aos eixos (prefix sum 2D, BIT, sweep).
  \item \textbf{Soma das distâncias entre todos os pares:} Manhattan já separa por eixo; para Chebyshev, use a volta e
        divida por $2$. Em cada eixo, ordene e some com prefix sums.
\end{itemize}
